% !TeX document-id = {20f24222-124f-4838-bced-b9cdf90be1c4}
% !TeX encoding = UTF-8
% !TeX spellcheck = de_DE

%%%%%%%%%%%%%%%%%%%%%%%%%%%%%%%%%%%%%%%%%%%%%%%%%%%%%%%%%%%%%%%%
% Beschreibung
%%%%%%%%%%%%%%%%%%%%%%%%%%%%%%%%%%%%%%%%%%%%%%%%%%%%%%%%%%%%%%%%
% Mit der folgenden Reihenfolge kompilieren, um alle Abhängigkeiten aufzulösen:
% 1) pdflatex
% 2) biber
% 3) pdflatex
% 4) pdflatex

% Die Datei "Bild-Metadaten-Analyzer.cls", "hsmwLogo.pdf" und "literatur.bib" müssen in einem Verzeichnis mit dieser Datei liegen.
% Ändern Sie die Datei "Bild-Metadaten-Analyzer.cls" bitte nicht, wenn nicht unbedingt notwendig!
% Die Datei "literatur.bib" beinhaltet Ihr Literaturverzeichnis und sollte dringend bearbeitet werden!
%%%%%%%%%%%%%%%%%%%%%%%%%%%%%%%%%%%%%%%%%%%%%%%%%%%%%%%%%%%%%%%%

\documentclass{Bild-Metadaten-Analyzer}

%%%%%%%%%%%%%%%%%%%%%%%%%%%%%%%%%%%%%%%%%%%%%%%%%%%%%%%%%%%%%%%%
% Konfigurationsbereich (wichtige Werte!)
% Bitte anpassen, sonst kompiliert das Dokument nicht!
%%%%%%%%%%%%%%%%%%%%%%%%%%%%%%%%%%%%%%%%%%%%%%%%%%%%%%%%%%%%%%%%

\Anrede{Frau}
\Vorname{Melissa}
\Nachname{Futtig}

\Matrikelnummer{107813}
\Seminargruppe{CC24w1-B}

% Betriebssysteme und digitale Spuren
% Netzwerkforensik
\Modulbezeichnung{Programmierung II} 

%%%%%%%%%%%%%%%%%%%%%%%%%%%%%%%%%%%%%%%%%%%%%%%%%%%%%%%%%%%%%%%%
% Literaturverwaltung
%%%%%%%%%%%%%%%%%%%%%%%%%%%%%%%%%%%%%%%%%%%%%%%%%%%%%%%%%%%%%%%%

% Verwendete Literatur bitte in diese Datei eintragen
\addbibresource{literatur.bib} % Name der Literatur-Datei
\usepackage[backend=biber,style=ieee]{biblatex}

%%%%%%%%%%%%%%%%%%%%%%%%%%%%%%%%%%%%%%%%%%%%%%%%%%%%%%%%%%%%%%%%
% Code-Formatierung
%%%%%%%%%%%%%%%%%%%%%%%%%%%%%%%%%%%%%%%%%%%%%%%%%%%%%%%%%%%%%%%%

\usepackage{listings}
\usepackage{xcolor}

\lstdefinestyle{pythonstyle}{
    language=Python,
    backgroundcolor=\color{gray!7},
    basicstyle=\ttfamily\small,
    keywordstyle=\color{blue},
    commentstyle=\color{green!50!black},
    stringstyle=\color{red!70!black},
    numbers=left,
    numberstyle=\tiny\color{gray},
    stepnumber=1,
    numbersep=8pt,
    frame=single,
    tabsize=4,
    breaklines=true
}

%%%%%%%%%%%%%%%%%%%%%%%%%%%%%%%%%%%%%%%%%%%%%%%%%%%%%%%%%%%%%%%%
% Dokument-Inhalt
%%%%%%%%%%%%%%%%%%%%%%%%%%%%%%%%%%%%%%%%%%%%%%%%%%%%%%%%%%%%%%%%

\begin{document}

% Kapitel 1
\section{Einleitung}\label{sec:Einleitung}
Die vorliegende Belegarbeit wird im Rahmen des Moduls ``Programmierung II`` des Studiengangs ``Bachelor of Science in IT-Forensik/ 
Cybercrime``, der Seminargruppe CC24w1-B, an der Hochschule Mittweida erstellt.\\
Dabei soll ein freigewähltes Programm in der Programmiersprache Python geschrieben werden, das sich erfolgreich implementieren und 
kompilieren lässt. Dieses trägt den Namen \textit{Bild-Metadaten-Analyzer}. In folgender \nameref{sec:Aufgabenstellung}
werden die Kriterien und Anforderungen an das Programm beschrieben. Im \nameref{sec:Ziel des Programms} aus Kapitel 
\Ref{sec:Ziel des Programms} werden die Programminhalte näher betrachtet.
\clearpage

% Kapitel 2
\section{Aufgabenstellung}\label{sec:Aufgabenstellung}
Das Programm sollte \ldots \\
\\\ldots in der Programmiersprache Python geschrieben sein.
\\\ldots mit einer Schnittstelle (GUI oder CLI) ausgestattet sein.
\\\ldots Inhalte verarbeiten und ausgeben können.
\\\ldots sich erfolgreich implementieren und kompilieren lassen.
\\\ldots plattformübergreifend laufend sein (Windows und Linux).
\\\ldots ohne kostenpflichtige Software oder Bibliothekne auskommen.
\\\ldots Quellcode mit beschreibenden Kommentaren enthalten.
\\\ldots gegebenenfalls Dateien enthalten, die zu Testzwecken genutzt werden können.\\
\\ Durch die freie Themenwahl gab es eine Vielzahl an Möglichkeiten, passende Programme zu erstellen. Damit dieses einen Bezug zum 
Studiengang hat, sollte es IT-forensische Inhalte beinhalten. So wurde sich für das Programm \textit{Bild-Metadaten-Analyzer} 
entschieden, das die Metadaten von Bildern analysieren und ausgeben kann. Es soll dabei helfen, Informationen über die Bilder zu 
erhalten, die für die IT-Forensik relevant sein können.
\clearpage

% Kapitel 3
\section{Ziel des Programms}\label{sec:Ziel des Programms}
Damit das Programm \textit{Bild-Metadaten-Analyzer} neben den Anforderungen aus der \nameref{sec:Aufgabenstellung} erfolgreich läuft,
sind weitere Kriterien festegelegt worden.
\\Unter anderem soll es \ldots\\
\\\ldots eine forensische Analyse von Bilddateien (*.jpeg, *.jpg, *.png, *.tif, und *.tiff) durchführen können.
\\\ldots die ausgelesenen Metadaten bzw. EXIF-Daten übersichtlich darstellen können.
\\\ldots über eine graphische Oberfläche (GUI), als benutzerfreundliche Schnittstelle, erreichbar sein.
\\\ldots die gefundenen Metadaten in einer \textit{*.csv-Datei} oder in einer \textit{*.json-Datei} extrahieren können.\\
\\EXIF-Daten sind Metadaten, die in Bilddateien enthalten sind. EXIF ist ein Akronym und steht für Exchangeable Image File Format.
``Dieses Dateiformat ist ein Standardformat und wird zum Speichern von Metadaten in digitalen Bildern verwendet. 
Diese Metadaten enthalten Informationen über das Bild, etwa die Einstellungen der Kamera wie Belichtungszeit, Blende und ISO-Wert. 
Aber auch Daten wie Datum, Uhrzeit und verwendete Kamera bei der Aufnahme. Wenn die Kamera mit einem GPS-System ausgestattet ist, 
können sogar die geografische Daten in den EXIF-Daten gespeichert sein.`` \cite{Onlinequelle:1}
\subsection{EXIF-Inhalte}\label{subsec:EXIF-Inhalte}
Wie in der Definition zu EXIF-Daten aus der Quelle \cite[E. M. Zylinski. „Was sind EXIF-Daten?“]{Onlinequelle:1}, werden in dem 
Programm \textit{Bild-Metadaten-Analyzer} die folgenden EXIF-Inhalte ausgelesen und dargestellt \ldots
\\\ldots Dateipfad mit Dateinamen,
\\\ldots Dateigröße,
\\\ldots Aufnahmedatum und -uhrzeit,
\\\ldots Kamerahersteller,
\\\ldots Kameramodell,
\\\ldots Objektiv,
\\\ldots Software,
\\\ldots Blende,
\\\ldots Belichtungszeit,
\\\ldots ISO-Wert,
\\\ldots Brennweite,
\\\ldots GPS-Daten.\\
\\Wenn EXIF-Daten nicht vorhanden sind, soll eine spezifische Fehlermeldung bzw. Information ausgegeben werden, sodass der Nutzer 
informiert ist, dass für das Bild keine EXIF-Daten vorliegen.

\subsection{Einsatzzweck}\label{subsec:Einsatzzweck}
Der sogenannte \textit{Bild-Metadaten-Analyzer} kann in der IT-Forensik von Nutzung sein, um Ermittlern der Polizei oder weiteren
nachrichtendienstlichen Behörden Informationen über Bilder zu liefern, die in einem Ermittlungsverfahren von Relevanz sein könnten.
\\Ein Anwendungsbeispiel hierbei wären Bilder, die einen Zeitstempel mit GPS-Daten enthalten. Solche könnten unterstützend wirken, 
um potentielle Täter oder Opfer eines Tatortes leichter zu identifizieren.
\clearpage

% Kapitel 4
\section{Anleitungen}\label{sec:Anleitungen}
In den folgenden Anleitungen (\nameref{subsec:Installationsanleitung}, \nameref{subsec:Konfigurationsanleitung} und 
\nameref{subsec:Bedienungsanleitung}) werden die Installation, Konfiguration und Bedienung des Programms 
\textit{Bild-Metadaten-Analyzer} beschrieben.

\subsection{Installationsanleitung}\label{subsec:Installationsanleitung}
Die Installationsanleitung enthält technische Voraussetzungen für die Installation auf einem MacOS, Windows oder Linux Betriebssystem. 
Es wird darauf eingegangen, wie das Programm zum ``Laufen'' gebracht wird.

\textbf{Technische Voraussetzungen:}
\begin{itemize}
	\item Python 3.10
	\item Betriebssystem: Windows 10/11, MacOS 10.15 oder höher, Linux (Ubuntu 20.04 oder höher)
	\item Mindestens 4 GB RAM
	\item Mindestens 100 MB freier Festplattenspeicher
\end{itemize}

\textit{Hinweis:} Das Programm wurde ausschließlich auf den oben genannten Betriebssystemen getestet und für die Entwicklung 
ein MacBook Pro mit MacOS 15.5 verwendet.

Um Python zu installieren, kann etweder die offizielle \href{https://www.python.org/downloads/}{Python-Webseite} besucht  
oder folgende Schritte aus Kapitel \ref{subsubsec:Installation von Bibliotheken} \nameref{subsubsec:Installation von Bibliotheken} 
durchlaufen werden.

\textbf{Benötigte Bibliotheken:}
\begin{itemize}
	\item Pillow (für die Bildverarbeitung bzw. die Bildübertragungsprozesse)
	\item Tkinter (für die GUI-Entwicklung und ist oft in der Python-Standardbibliothek enthalten)
\end{itemize}

\textbf{Standardbibliotheken}\\
Folgende Standardbibliotheken werden für die Entwicklung für das Programm \textit{Bild-Metadaten-Analyzer} verwendet, kommen 
jedoch mit der Installation von Python (siehe Kapitel \ref{subsubsec:Installation von Bibliotheken} \nameref{subsubsec:Installation von Bibliotheken})einher:
\begin{itemize}
	\item JSON: für die Verarbeitung von JSON-Daten --> zum Speichern von Dateien im JSON-Format
	\item CSV: für die Verarbeitung von CSV-Daten --> zum Speichern von Dateien im CSV-Format
\end{itemize}
\subsubsection{Installation von Bibliotheken}\label{subsubsec:Installation von Bibliotheken}

\textbf{Installation auf Windows:}\\
\\\textit{Python}\\
Für die Installation von Python auf einem Windows-Betriebssystem wurde die Dokumentation von 
\href{https://docs.python.org/3.10//using/windows.html}{Python} selbst genutzt. Die einzelnen Schritte werden hier nicht näher 
beschrieben, um den Rahmen der Belegarbeit nicht auszureizen.

\textit{TKinter}\\
Die Installation von TKinter unter Windows erfordert Python 3.10 oder höher als Vorinstallation. Diese Bibliothek lässt sich über 
die Windows Eingabeaufforderung (CMD) mit folgendem Befehl installatieren:
\begin{lstlisting}[style=pythonstyle]
pip install python-tk
\end{lstlisting}

\textit{Pillow}\\
Die Installation von Pillow unter Windows erfordert Python 3.10 oder höher als Vorinstallation. Diese Bibliothek lässt sich über 
die Windows Eingabeaufforderung (CMD) mit folgendem Befehl installatieren:
\begin{lstlisting}[style=pythonstyle]
pip install pillow
\end{lstlisting}

\textit{Hinweis:} Für den Fall, dass `pip' nicht erkannt wird, kann auch `pip3' verwendet werden.\\

\textbf{Installation auf Linux:}\\
\\\textit{Python}\\
Für die Installation von Python auf einem Linux-Betriebssystem wurde sich an der Dokumentation von 
\href{https://computingforgeeks.com/how-to-install-python-on-ubuntu-linux-system/}{Computing for Geeks} für ein Ubuntu 20.04 und 
höher orientiert. Die Installation erfolgt nicht über die GUI, sondern die Kommandozeile (Terminal).\\
\\Mit folgenden Befehelen kann Python installiert werden:

\begin{lstlisting}[style=pythonstyle]
sudo apt install software-properties-common -y

sudo add-apt-repository ppa:deadsnakes/ppa

sudo apt install python3.10
\end{lstlisting}

Mit folgendem Befehl erfolgt die Überprüfung, ob Python erfolgreich installiert wurde:
\begin{lstlisting}[style=pythonstyle]
python3.10 --version
\end{lstlisting}
\textit{Hinweis:} Wenn Befehle mit `sudo' ausgeführt werden, sollte das Passwort des Benutzers bereitgehalten und eingegeben werden.\\

\textit{TKinter}\\
Die Installation von TKinter unter Linux erfordert Python 3.10 oder höher als Vorinstallation. Diese Bibliothek lässt sich über 
das Linux Terminal mit folgendem Befehl installatieren:
\begin{lstlisting}[style=pythonstyle]
sudo apt install python3-tk
\end{lstlisting}

\textit{Pillow}\\
Die Installation von Pillow unter Linux erfordert Python 3.10 oder höher als Vorinstallation. Diese Bibliothek lässt sich über 
das Linux Terminal mit folgendem Befehl installatieren:
\begin{lstlisting}[style=pythonstyle]
sudo apt install pillow
\end{lstlisting}
oder auch \\
\begin{lstlisting}[style=pythonstyle]
pip install pillow
\end{lstlisting}

\subsection{Konfigurationsanleitung}\label{subsec:Konfigurationsanleitung}
Die Konfigurationsanleitung beschreibt welche Einstellungen als Nutzer selbst vorgenommen werden könnten. Als mögliche Einstellungen 
werden unter anderem die Bildformate betrachtet, die ausschließlich analysiert werden können. Dabei handelt es sich um die Formate, 
welche bereits in Kapitel \ref{sec:Ziel des Programms} \nameref{sec:Ziel des Programms} gelistet wurden.\\

Des Weiteren hat der Nutzer die Wahl, ob die ausgelesenen Metadaten als \textit{*.csv-Datei} oder als \textit{*.json-Datei} 
extrahiert werden sollen. Der Speicherort kann hierbei \textbf{nicht} frei gewählt werden und ist im Programm-Code festgelegt und 
befindet sich \textbf{immer} im DOWNLOADS-Ordner eines jeden Systems.\\

Zusätzlich ist die Namensgebung der extrahierten Dateien festgelegt. Die Dateien behalten ihren ursprünglichen Dateinamen, jedoch 
mit der Endung \textit{*\_Metadata\_Auswertung.json} oder \textit{*\_Metadata\_Auswertung.csv}, je nachdem, für welches Format sich der 
Nutzer entscheidet.\\

Die Ausgabe der Sprache lässt sich nicht verändern. Die Originalsprache des Systems der Bilder ist englisch und wird in das Deutsche 
übersetzt.

\subsection{Bedienungsanleitung}\label{subsec:Bedienungsanleitung}
In wenigen Schritten ist es einem Endnutzer diesen Programmes möglich dieses auszuführen.\\
\\\textbf{Schritt 1:}
\\Das Programm \textit{Bild-Metadaten-Analyzer} wird über die ausführbare Datei \textit{main.py} gestartet. Ein Weg wäre es  
über die Eingabeaufforderung (CMD) unter Windows oder über das Terminal unter Linux auszuführen ODER über eine Intelligente 
Entwicklungsumgebung (IDE) wie zum Beispiel IntelliJ oder Visual Studio Code.\\
\\Gestartet wird das Programm mit folgendem Befehl:
\begin{lstlisting}[style=pythonstyle]
python main.py
\end{lstlisting}
\textit{Alternativ auch mit:}
\begin{lstlisting}[style=pythonstyle]
python3 main.py
\end{lstlisting}

\textbf{Schritt 2:}\\
Nach dem Starten des Programms öffnet sich die graphische Oberfläche (GUI) des Programms. Es erscheint ein Fenster, in dem der 
Endnutzer auswählen kann, welche Bilddatei betrachtet werden soll. Alle Dateien, die nicht kompatibel sind, werden automatisch 
ausgegraut, sodass eine Auswahl dieser nicht erfolgen kann.
\section{Programmaufbau}

\subsection{Vorstellung Quellcode-Dateien}\label{subsec:Vorstellung Quellcode-Dateien}
dummytext

\subsubsection{Software-Architektur}\label{subsubsec:Software-Architektur}

\subsection{Funktionsumfang}\label{subsec:Funktionsumfang}

\begin{itemize}
	\item 
	\item 
	\item
\end{itemize}

\subsection{Programmablauf}\label{subsec:Programmablauf}

dummytext

\section{Fazit}\label{sec:Fazit}
dummytext

\subsection{Verwendete Hilfsmittel}\label{subsec:Verwendete Hilfsmittel}

\subsection{Technische Herausforderungen}\label{subsec:Technische Herausforderungen}

\subsection{Weitere Features}\label{subsec:WeitereFeatures}
Überprüfung, ob die EXIF-Inhalte verändert worden sind

\end{document}